\documentclass[../main.tex]{subfiles}
\begin{document}

\chapter*{RESUMEN}

Este trabajo surge en el contexto de la creciente implementación del propano (R290) en el mercado de las bombas de calor, impulsada por la entrada en vigor del Reglamento F-Gas (UE 2024/573), que establece límites estrictos a la utilización de HFC con elevado potencial de calentamiento global.

Si bien el R290 ha ganado protagonismo gracias a su buen rendimiento energético y su bajo GWP, existe margen para mejorar sus prestaciones mediante el uso de mezclas de refrigerantes naturales. En este trabajo se lleva a cabo un análisis teórico y experimental de diversas mezclas alternativas al R290 orientadas a la producción de calor.

Este trabajo demuestra experimentalmente la viabilidad de las mezclas zeotrópicas de dimetil éter y propileno en un rango de composiciones que va desde 95\%/5\% hasta 80\%/20\% en base másica, siendo esta última composición (80\%/20\%) la que ha proporcionado los mejores resultados. Con dicha mezcla se ha alcanzado experimentalmente un incremento del rendimiento energético (COP) de hasta un +16\% respecto al R290 en el rango de temperaturas de agua de 55°C a 65°C, y de hasta un +12\% en el rango de 47ºC a 55ºC. Estos resultados vienen acompañados de una reducción de la capacidad calorífica volumétrica (VHC) del -22\% y -23\% respecto al propano, respectivamente, para las temperaturas mencionadas.

\chapter*{RESUM}

Este treball sorgeix en el context de la creixent implementació del propà (R290) en el mercat de les bombes de calor, impulsada per l'entrada en vigor del Reglament F-Gas (UE 2024/573), que estableix límits estrictes a la utilització d'HFC amb elevat potencial d'escalfament global.

Si bé el R290 ha guanyat protagonisme gràcies al seu bon rendiment energètic i al seu baix GWP, hi ha marge per a millorar les seues prestacions mitjançant l'ús de mescles de refrigerants naturals. En este treball es duu a terme un anàlisi teòric i experimental de diverses mescles alternatives al R290 orientades a la producció de calor.

Este treball demostra experimentalment la viabilitat de les mescles zeotròpiques de dimetil èter i propilé en un rang de composicions que va des de 95\%/5\% fins a 80\%/20\% en base màssica, sent esta última composició (80\%/20\%) la que ha proporcionat els millors resultats. Amb la dita mescla s'ha aconseguit experimentalment un increment del rendiment energètic (COP) de fins a un +16\% respecte al R290 en el rang de temperatures d'aigua de 55°C a 65°C, i de fins a un +12\% en el rang de 47ºC a 55ºC. Estos resultats vénen acompanyats d'una reducció de la capacitat calorífica volumètrica (VHC) del -22\% i -23\% respecte al propà, respectivament, per a les temperatures mencionades.

\chapter*{ABSTRACT}

This work arises in the context of the growing implementation of propane (R290) in the heat pump market, driven by the entry into force of the F-Gas Regulation (EU 2024/573), which establishes strict limits on the use of HFCs with high global warming potential.

Although R290 has gained prominence due to its good energy performance and low GWP, there is room for improvement in its performance through the use of natural refrigerant mixtures. This work carries out a theoretical and experimental analysis of various alternative mixtures to R290 aimed at heat production.

This work experimentally demonstrates the viability of zeotropic mixtures of dimethyl ether and propylene across a range of compositions from 95\%/5\% to 80\%/20\% by mass, with the latter composition (80\%/20\%) providing the best results. With this mixture, an experimental increase in energy performance (COP) of up to +16\% compared to R290 has been achieved in the water temperature range of 55°C to 65°C, and up to +12\% in the range of 47ºC to 55ºC. These results are accompanied by a reduction in volumetric heating capacity (VHC) of -22\% and -23\% compared to propane, respectively, for the temperatures mentioned.


\newpage
\end{document}